\section{Related Work}
\label{sec:related}

Formalizing statutory rules as executable logic dates to Sergot et al.'s encoding of the British Nationality Act as a Prolog program~\cite{sergot1986}, and continues through systems such as PROLEG~\cite{satoh2010proleg} and defeasible-argumentation frameworks~\cite{benchcapon2007,prakken2020}, whose soundness and explainability depend on a correctly and completely specified rule structure - traditionally hand-entered rather than extracted from unstructured text. The rise of LLMs has shifted attention toward automatically bridging this gap. LegalBench~\cite{guha2023legalbench} evaluates LLMs on legal reasoning tasks via direct or chain-of-thought prompting for a final label, an evaluation paradigm whose reliability is questionable: chain-of-thought explanations can be systematically unfaithful to the process that actually produced a model's answer~\cite{turpin2023unfaithful}, and a related faithfulness study of LLMs as autoformalizers in legal reasoning documents an analogous divergence between benchmark accuracy and logical faithfulness~\cite{wang2026knowyourlimits}. Closer to this paper's architecture, several systems couple LLMs with an explicit symbolic component rather than relying on the LLM's final-answer judgment alone, e.g., knowledge-graph-guided legal provision prediction~\cite{li2021legalkg} and a multi-agent framework that formalizes statutes into verifiable intermediate representations before symbolic inference~\cite{sadowski2025verifiable}. These works establish that routing LLM output through a symbolic intermediate is broadly beneficial, consistent with what this paper also finds, but they evaluate only whether the final decision is correct, not which intermediate extraction errors were responsible when it is not - the question this paper addresses through predicate-level intervention on a fixed, auditable rule tree.
